\section{Πρόβλημα 7}

\subsection{Ερώτηση (i)}

Υλοποιήσαμε τον αλγόριθμο Fictitious Play σε Python. Κάθε παίκτης διατηρεί ένα
ιστορικό των ενεργειών του αντιπάλου και σε κάθε γύρο $t$ παίζει την βέλτιστη
απόκριση στην εμπειρική κατανομή του αντιπάλου μέχρι εκείνη τη στιγμή. Ο
αλγόριθμος τερματίζει είτε όταν φτάσουμε $T_{\max}$ γύρους, είτε όταν και οι δύο
παίκτες έχουν regret μικρότερο ή ίσο του $\varepsilon$, δηλαδή:
$$
\max_i (A\mathbf{q})_i - \mathbf{p}^\top A \mathbf{q} \leq \varepsilon
\quad \text{και} \quad
\max_j (B^\top \mathbf{p})_j - \mathbf{q}^\top B^\top \mathbf{p} \leq \varepsilon,
$$
όπου $\mathbf{p}, \mathbf{q}$ είναι οι εμπειρικές κατανομές των παικτών 1 και 2
αντίστοιχα, $A$ ο πίνακας ωφελειών του παίκτη 1 (γραμμές = στρατηγικές παίκτη 1,
στήλες = στρατηγικές παίκτη 2) και $B^\top$ ο αντίστοιχος πίνακας για τον παίκτη
2 (γραμμές = στρατηγικές παίκτη 2, στήλες = στρατηγικές παίκτη 1). Στη δεύτερη
περίπτωση, το παίγνιο έχει συγκλίνει σε $\varepsilon$-σημείο ισορροπίας.
Ορίσαμε $T_{\max} = \min(2^n, 10^6)$.

\subsection{Ερώτηση (ii)}

Για κάθε τιμή $n \in \{20, 50, 100\}$ παράγαμε 10 τυχαία $n \times n$ παίγνια
μηδενικού αθροίσματος με $B = -A$ και $A_{ij} \in [0,1]$, και τρέξαμε τον
αλγόριθμο για $\varepsilon \in \{0.01, 0.001\}$.

\begin{center}
\begin{tabular}{cccccc}
    $n$ & $\varepsilon$ & $T_{\max}$ & Σύγκλιση & Μέσος αρ. επαναλήψεων & Μέσος αρ. (σύγκλιση) \\
    \hline
    20  & 0.01  & 1\,000\,000 & 10/10 & 3\,521  & 3\,521   \\
    20  & 0.001 & 1\,000\,000 & 10/10 & 466\,851 & 466\,851 \\
    50  & 0.01  & 1\,000\,000 & 10/10 & 7\,758  & 7\,758   \\
    50  & 0.001 & 1\,000\,000 & 2/10  & 980\,607 & 903\,034 \\
    100 & 0.01  & 1\,000\,000 & 10/10 & 15\,921 & 15\,921  \\
    100 & 0.001 & 1\,000\,000 & 0/10  & 1\,000\,000 & --- \\
\end{tabular}
\end{center}

Για $\varepsilon = 0.01$ ο αλγόριθμος συγκλίνει πάντα και για τις τρεις τιμές
του $n$, επιβεβαιώνοντας τη θεωρητική εγγύηση σύγκλισης για παίγνια μηδενικού
αθροίσματος. Ο μέσος αριθμός επαναλήψεων κλιμακώνεται περίπου γραμμικά με το
$n$ (3\,521, 7\,758, 15\,921), υποδηλώνοντας αργή αλλά σταθερή σύγκλιση. Για
$\varepsilon = 0.001$, το $T_{\max} = 10^6$ δεν επαρκεί για $n \geq 50$: η
μείωση του $\varepsilon$ κατά παράγοντα 10 απαιτεί πολύ περισσότερους γύρους,
υποδηλώνοντας ότι ο αλγόριθμος συγκλίνει αργά ως προς το $\varepsilon$.

\subsection{Ερώτηση (iii)}

Επαναλάβαμε το πείραμα για τυχαία $n \times n$ παίγνια γενικού αθροίσματος
($A, B$ ανεξάρτητα με $A_{ij}, B_{ij} \in [0,1]$).

\begin{center}
\begin{tabular}{cccccc}
    $n$ & $\varepsilon$ & $T_{\max}$ & Σύγκλιση & Μέσος αρ. επαναλήψεων & Μέσος αρ. (σύγκλιση) \\
    \hline
    20  & 0.01  & 1\,000\,000 & 9/10  & 108\,172 & 9\,080    \\
    20  & 0.001 & 1\,000\,000 & 8/10  & 334\,898 & 168\,623  \\
    50  & 0.01  & 1\,000\,000 & 5/10  & 501\,492 & 2\,984    \\
    50  & 0.001 & 1\,000\,000 & 8/10  & 304\,721 & 130\,901  \\
    100 & 0.01  & 1\,000\,000 & 7/10  & 354\,838 & 78\,340   \\
    100 & 0.001 & 1\,000\,000 & 2/10  & 853\,389 & 266\,943  \\
\end{tabular}
\end{center}

Σε αντίθεση με τα παίγνια μηδενικού αθροίσματος, η σύγκλιση δεν είναι
εγγυημένη θεωρητικά για γενικά παίγνια, και αυτό επιβεβαιώνεται εμπειρικά.
Παρατηρούμε τα εξής:

\begin{itemize}

    \item Σε αντίθεση με τα παίγνια μηδενικού αθροίσματος όπου η σύγκλιση
        είναι εγγυημένη, εδώ ορισμένα παίγνια δεν συγκλίνουν εντός
        $T_{\max}$: π.χ.\ για $n=50$, $\varepsilon=0.01$ μόνο 5/10 παίγνια
        συγκλίνουν, και για $n=100$, $\varepsilon=0.001$ μόνο 2/10. Αυτό
        είναι αναμενόμενο καθώς δεν υπάρχει θεωρητική εγγύηση σύγκλισης για
        γενικά παίγνια.

    \item Για τα παίγνια που συγκλίνουν, ο αριθμός επαναλήψεων αυξάνεται
        δραματικά καθώς μικραίνει το $\varepsilon$: για $n=20$, ο μέσος
        αριθμός επαναλήψεων πηγαίνει από 9\,080 για $\varepsilon=0.01$ σε
        168\,623 για $\varepsilon=0.001$, δηλαδή αύξηση κατά $\sim\!18.6\times$
        για μείωση του $\varepsilon$ κατά $10\times$. Η ίδια τάση παρατηρείται
        και στα παίγνια μηδενικού αθροίσματος (3\,521 έναντι 466\,851 για
        $n=20$, αύξηση $\sim\!133\times$), γεγονός που υποδηλώνει αργή
        σύγκλιση ως προς το $\varepsilon$ ανεξαρτήτως τύπου παιγνίου.

\end{itemize}
